\section{Practical Recommendations}
\label{sec:discussion}

Production deployment requires considering
operational properties alongside statistical fidelity. Historical caching
retains the full question set, but requires recent outcomes and monitoring for stale verdicts. Fixed subsets provide predictable
execution and common questions across runs. Adaptive testing can
tailor selection to each run, but requires sequential orchestration,
run-specific state, and repeated ability updates. These operational differences
make the appropriate method dependent on the surrounding evaluation workflow.


For our production workflow, we deployed difficulty-stratified fixed subsets with $k\in\{100,200,300,400\}$, allowing users to choose an execution--fidelity tradeoff. 
Fixed subsets reveal the workload
before execution and preserve question-level comparisons for diagnosing
behavioral changes. They also avoid additional sequential serving logic in an evaluation
system designed to execute questions in parallel.
The subset transferred without recalibration to five additional agent families
and showed no monotonic loss of fidelity across the calibration windows. These
results support reuse across the systems and period studied.

As the agent evolves, we recommend recalibrating after material changes to its
models, prompts, tools, or execution system, or after sustained shifts in
question-level outcomes. Developers should also periodically run the full
benchmark and compare its pass rate with the fixed-subset estimate to measure
the error introduced during ongoing monitoring.